\subsection{All sections of the intrinsic Heisenberg quotient}

\begin{statement}{E.35}{Proved}{A norm equation and every conjugacy class in the finite intrinsic quotient}
Let $K=\mathbb Q(\zeta)$ and retain the quotient $H_3$ with
the based $G_K$-action (E31.3).
Write $\mathcal A=K^\times/K^{\times3}$ multiplicatively.
For an abelian section pair $(a,b)\in\mathcal A^2$ put
$A=\zeta a$, choose nonzero representatives $A,b\in K$,
and form the etale cubic algebra
$L_A=K[t]/(t^3-A)$.
The fiber of section conjugacy classes of
$H_3\rtimes G_K$ over $(a,b)$ is nonempty exactly when
the explicit equation
\[
 X^3+A Y^3+A^2 Z^3-3AXYZ=b
                 \quad\text{has a solution }X,Y,Z\in K .
 \tag{E35.1}
\]
Every nonempty fiber is a torsor under the explicitly
specified group
\[
 \mathcal A/\langle A,b\rangle
       =K^\times/(K^{\times3}\langle A,b\rangle).
 \tag{E35.2}
\]
The subgroup in the denominator has rank zero, one or two
over $\mathbb F_3$, according to the span of the two
indicated cube classes.
There are concrete coordinates for this group. Choose one
generator $\pi_{\mathfrak p}$ for each nonzero prime ideal of
$\mathbb Z[\zeta]$, and put
$J=\{\zeta\}\amalg\{\mathfrak p\}$.
Unique factorization gives the explicit isomorphism
\[
 \mathcal A\longrightarrow\bigoplus_{j\in J}\mathbb F_3,
 \qquad
 [\zeta^e\prod_{\mathfrak p}\pi_{\mathfrak p}^{e_{\mathfrak p}}]
       \longmapsto(e,(e_{\mathfrak p}))\pmod3 .
 \tag{E35.2a}
\]
The sign in the factorization is a cube and is omitted.
Thus (E35.2) is obtained by Gaussian elimination of the
two finite-support vectors $[A],[b]$: choose zero, one or
two pivot coordinates according to their rank, subtract
the corresponding row-reduced basis vectors, and retain
all remaining coordinates. This is an explicit
parameterization by the direct sum of $\mathbb F_3$ over
the coordinates in $J$ outside those pivots.
After choosing the lift $n_0$ constructed from a norm solution,
all classes are represented by
$n_d(g)=(0,0,\kappa_3(d)(g))n_0(g)$ with
$d\in K^\times/(K^{\times3}\langle A,b\rangle)$.

All dependent cases are included. If $A$ is a cube,
(E35.1) is always soluble and the fiber group is
$\mathcal A/\langle b\rangle$.
If $A$ is not a cube and $b$ is in its span, (E35.1) is
always soluble and the fiber group is
$\mathcal A/\langle A\rangle$.
If both classes are trivial the group is $\mathcal A$.
For two independent classes, a solution of (E35.1)
constructs an actual Heisenberg Galois extension as in the proof.

The quotient is the one obtained from the actual map (E31.4)
and the based arithmetic comparison of E.31.
For example the pair $(a,b)=(1,7)$ has empty fiber,
while $(a,b)=(2\zeta^2,3)$ has nonempty fiber, with
norm solution $1+\sqrt[3]{2}$.
No assertion is made that every pair here is the image of
a class in the full intrinsic rank-ten abelian section group.

\emph{Manuscript: \S51.16, folios 733v--736v, scans 172--178,
(371)--(379); this computes a specified finite intrinsic
quotient reached by the additional morphism (E31.4).}
\end{statement}
\begin{proof}
Choose the path in E.31 within the formal punctured disc
at the target zero. Equation (E31.7) then gives the image
of the source tangential section as
\[
 q_g=(\kappa_3(\zeta)(g),0,0)\in H_3 .
 \tag{E35.3}
\]
To see the zero central coordinate directly, write
$h=\zeta\lambda^{-3}s^3u(s)$ with $u(0)=1$.
Its chosen cube root is
$(\zeta_9/\lambda)s\,u(s)^{1/3}$.
The last factor has coefficients in $K$, so the only
Galois change is the inertia power supplied by
$\zeta_9$. The formal puncture has only that one
geometric generator and gives no commutator term.
Thus $q:G_K\to H_3$ is the homomorphism in (E35.3),
and the intrinsic action is $\operatorname{Ad}(q_g)$.

There is an explicit isomorphism of arithmetic groups
\[
 H_3\rtimes_{\operatorname{Ad}q}G_K
       \longrightarrow H_3\times G_K,\qquad
 (n,g)\longmapsto(nq_g,g).
 \tag{E35.4}
\]
Its inverse sends $(h,g)$ to $(hq_g^{-1},g)$.
It is the identity on the geometric group.
Consequently it identifies section classes with continuous
homomorphisms $G_K\to H_3$ modulo conjugation by $H_3$,
and sends the intrinsic abelian pair $(a,b)$ to the
ordinary Kummer pair $(A,b)=(\zeta a,b)$.

Write these latter Kummer characters as $u,v:G_K\to\mathbb F_3$.
A homomorphism with those abelian coordinates has the form
$(u_g,v_g,w_g)$, where its central coordinate solves
\[
 d w(g,h)=-\tfrac12(u_gv_h-v_gu_h).
 \tag{E35.5}
\]
The two-cocycle on the right, with sign reversed, differs
from $u_gv_h$ by
$d(u_gv_g/2)$; hence its obstruction is exactly
$\kappa_3(A)\cup\kappa_3(b)$.
The Kummer exact sequence and Hilbert 90 embed this group
of classes into $\operatorname{Br}(K)[3]$.

We spell out the resulting norm condition.
Suppose first that $A$ is not a cube, let
$\alpha^3=A$ and let $\sigma(\alpha)=\zeta\alpha$ generate
$\operatorname{Gal}(L_A/K)$.
The cyclic algebra is
\[
 \mathcal C=L_A\oplus L_A U\oplus L_A U^2,\qquad
 U\ell=\sigma(\ell)U,\quad U^3=b .
 \tag{E35.6}
\]
Its factor set is
$b^{\lfloor(i+j)/3\rfloor}$ on $\sigma^i,\sigma^j$.
Choose $\beta^3=b$ and represent the values of $u$ by
$0,1,2$. The multiplicative coboundary of
$g\mapsto\beta^{u_g}$ is
\[
 b^{\lfloor(u_g+u_h)/3\rfloor}\zeta^{v_gu_h}.
 \tag{E35.7}
\]
The identity for $d(u_gv_g)$ in (E35.5) shows that the
class of this factor set is precisely
$\kappa_3(A)\cup\kappa_3(b)$.
The crossed-product comparison with the Brauer group is
Milne, \emph{Class Field Theory}, IV.3.14
\cite{E_MilneCFT}; its construction uses exactly the
multiplication and factor set just displayed.

The algebra (E35.6) is split exactly when $b$ is a norm.
This equivalence also has a direct proof here.
A norm solution $b=N_{L_A/K}(w)$ represents $\mathcal C$
on the three-dimensional $K$-space $L_A$, with $\ell$
acting by multiplication and $U$ acting as $w\sigma$.
The relations hold and the resulting homomorphism is an
isomorphism with $\operatorname{End}_K(L_A)$, by central
simplicity and the equal dimension nine.
Conversely, a splitting realizes $\mathcal C$ on a
three-dimensional $K$-space which, by its embedded
degree-three field $L_A$, is one-dimensional over $L_A$.
The operator $U$ is necessarily $w\sigma$ in a choice
of $L_A$-basis. Its cube is $N(w)$, so $b=N(w)$.
The norm of $X+Y\alpha+Z\alpha^2$, by its multiplication
matrix on this same basis, is
$X^3+A Y^3+A^2Z^3-3AXYZ$.
This proves (E35.1).

For independent cube classes $A,b$ the norm solution
also yields an explicit finite Galois realization.
Write $w=X+Y\alpha+Z\alpha^2$ with $N(w)=b$ and put
\[
 \beta^3=b,\qquad
 v=w^2\sigma(w),\qquad \gamma^3=v.
 \tag{E35.8}
\]
The field $M=L_A(\beta,\gamma)$ has the automorphisms
\[
 \begin{array}{c|ccc}
       &\alpha&\beta&\gamma\\ \hline
 \widetilde\sigma&\zeta\alpha&\beta&\beta\gamma/w\\
 \tau&\alpha&\zeta\beta&\gamma .
 \end{array}
 \tag{E35.9}
\]
Indeed $\sigma(v)/v=b/w^3$, so these formulas preserve
the defining equations, and
$\widetilde\sigma^3(\gamma)
=b\gamma/N(w)=\gamma$.
Both displayed automorphisms have order three.
The root $\gamma$ does not already lie in $L_A(\beta)$.
Otherwise the fact that its cube is in $L_A$ would write
$\gamma=c\beta^j$ for $c\in L_A$.
The equation $\sigma(v)/v=b/w^3$ would then imply that
$b$ is a cube in $L_A$, contradicting the independence
of $A,b$.
Thus $[M:K]=27$, the formulas make $M/K$ Galois, and
their commutator acts on $\gamma$ as multiplication by
$\zeta^{-1}$. Taking that automorphism as the central
generator identifies its Galois group with (E31.2),
with the two stated Kummer characters. Composing the
Galois restriction map with this identification and then
(E35.4)$^{-1}$ gives an actual arithmetic section.

The dependent cases require no degree-27 field.
If $A$ is a cube, replace its representative by one.
Then $L_A\simeq K^3$ and the norm is the product of its
three coordinates, so it is surjective. In (E35.1) the
choice
$X=(b+2)/3$, $Y=Z=(b-1)/3$ gives norm $b$.
The homomorphism $g\mapsto(0,\kappa_3(b)(g),0)$
provides a lift.
If $A$ is nontrivial and $b=A^j c^3$ for $j=0,1,2$,
the element $w=c\alpha^j$ has norm $b$.
The homomorphism
$g\mapsto(\kappa_3(A)(g),j\kappa_3(A)(g),0)$
has the required coordinates; its alternating term is
zero. This includes $b=1$ and all other dependent
cases. The construction is independent of the selected
representatives modulo cubes.

Finally choose one lift $h_g$ of $(u,v)$.
Every other lift is uniquely $z_g h_g$ with
$z\in H^1(G_K,\mathbb F_3)=\mathcal A$.
Conjugating by an element with abelian coordinates
$(r,s)$ adds the central character
$r v-s u$. Central conjugations do nothing.
The exact subgroup of differences generated by all
geometric conjugations is therefore
$\langle\kappa_3(A),\kappa_3(b)\rangle$.
This proves the entire torsor description (E35.2),
including each of the dependent cases just treated.
For completeness, the coordinate computation (E35.2a)
uses the Euclidean norm
$N(a+b\zeta)=a^2-ab+b^2$ on $\mathbb Z[\zeta]$.
Rounding the two rational coefficients of any quotient
to nearest integers leaves errors $r,s$ of absolute
value at most $1/2$, and
$r^2-rs+s^2\leq3/4<1$. The division algorithm proves
that this ring is a principal ideal domain and a unique
factorization domain. Its norm-one units are
$\{\pm1,\pm\zeta,\pm\zeta^2\}$, so their quotient by
cubes is generated freely over $\mathbb F_3$ by $\zeta$.
Factoring the numerator and denominator of an element
of $K^\times$ now proves both the injectivity and
surjectivity of (E35.2a). Row reduction gives an actual
isomorphism from its quotient by $[A],[b]$ to the stated
direct sum: normalize the pivot entries to one and all
other pivot entries to zero, then send a vector to its
unique representative with every pivot entry zero.
This also describes the inverse by inclusion of those
nonpivot coordinates.

At $(a,b)=(1,7)$, equation (E35.1) would assert that
$7$ is a norm from $K(\zeta_9)$, contradicted by the
local calculation at $(3+\zeta)$ in E.31.
For $A=2,b=3$, the element $1+\alpha$ has norm
$1+\alpha^3=3$, proving the positive example.
The program \texttt{e\_intrinsic\_finite.py} checks the
norm polynomial, this solution and the relations in
(E35.8)--(E35.9), and enumerates every conjugation
orbit for all two-generator assignments in $H_3$.
The finite intrinsic quotient map of E.31 remains the
comparison to the full cubic problem; no inverse to
that quotient map is asserted.
\end{proof}
